\documentclass[11pt]{article}
\usepackage[margin=0.9in]{geometry}
\usepackage[T1]{fontenc}
\usepackage{lmodern}
\usepackage{microtype}
\usepackage{booktabs}
\usepackage{enumitem}
\usepackage{xcolor}
\usepackage[hidelinks]{hyperref}

\definecolor{eyeblue}{HTML}{203B78}
\definecolor{eyegreen}{HTML}{1B756B}
\setlist{nosep,leftmargin=*}
\setlength{\parskip}{0.45em}
\setlength{\parindent}{0pt}

\newcommand{\choices}[5]{\begin{enumerate}[label=\Alph*.,leftmargin=2em,itemsep=0.15em,topsep=0.2em]\item #2\item #3\item #4\item #5\end{enumerate}}
\newcommand{\tasktitle}[1]{\vspace{0.4em}\noindent\textbf{#1}\par}
\newcommand{\fillline}[1]{\rule{#1}{0.4pt}}

\title{\textcolor{eyeblue}{EyeExplore User Study}\\
\large Four-Set Multiple-Choice Task Booklet}
\author{}
\date{}

\begin{document}
\maketitle

\section*{Participant instructions}
You will complete four short questions using the EyeExplore dashboard. Follow
the steps given by the study leader. Each question has one correct answer.
Choose exactly one option for every question. You do not need to write an
explanation.

The study leader will set up the dataset, stimulus, records, and dashboard
views before you begin. Do not change those settings unless the leader asks
you to. If a required view does not load, tell the study leader and wait.
An AOI is a marked region on the image. A visit means looking within that
region continuously before moving away. ``Focal'' means focused inspection;
``ambient'' means broader search.

\section*{Study leader: use before the session}
\textbf{Do not give this section or the answer key to participants.} Each
participant completes exactly one of the four sets. Randomly assign novices
between Sets 1 and 2, and experts between Sets 3 and 4. Give every participant
the same five-minute introduction. Keep the study data and settings fixed
within matched sets.

\subsection*{Set-up card - fill and freeze before recruiting}
Use symbolic record labels in the booklet and map them to real IDs on the
private operator copy. Use the same mapping for all participants.

\begin{center}
\begin{tabular}{p{0.34\linewidth}p{0.55\linewidth}}
\toprule
Study item & Locked value \\
\midrule
Dataset and version & \fillline{8cm} \\
Image stimulus & \fillline{8cm} \\
Target AOI name and boundary & \fillline{8cm} \\
Comparison AOI name and boundary & \fillline{8cm} \\
Video stimulus / episode A & \fillline{8cm} \\
Video stimulus / episode B & \fillline{8cm} \\
Video stimulus / episode C & \fillline{8cm} \\
Video stimulus / episode D & \fillline{8cm} \\
Record P1 actual ID & \fillline{8cm} \\
Record P2 actual ID & \fillline{8cm} \\
Record P3 actual ID & \fillline{8cm} \\
Record P4 actual ID & \fillline{8cm} \\
\bottomrule
\end{tabular}
\end{center}

\textbf{Definitions used for scoring}
\begin{itemize}
  \item A fixation is assigned to an AOI when its gaze point falls inside the
  locked AOI boundary. Use this rule for every participant record.
  \item A visit to an AOI is one uninterrupted run of fixations in that AOI.
  Consecutive fixations in the same AOI count as one visit. Leaving and later
  returning counts as a new visit.
  \item Total dwell time in an AOI is the sum of the durations of all fixations
  assigned to that AOI. Do not include time outside the AOI.
  \item Dwell share is total dwell time in the Target AOI divided by total
  fixation duration for that record, multiplied by 100.
  \item A directed transition from AOI X to AOI Y occurs when one visit to X is
  followed immediately by a visit to Y. Count each X-to-Y change once. Do not
  count X-to-X as a transition.
  \item For video episodes, a fixation is focal when its K-Coefficient is above
  zero, ambient when it is below zero, and neutral when it equals zero. An
  episode is primarily focal only when more than 50\% of its total fixation
  duration has a positive K-Coefficient. Use fixation duration as the weight.
\end{itemize}

\textbf{Layout conditions}
\begin{itemize}
  \item \textbf{One-view condition:} show only one visualization at a time.
  Participants may switch views when the task instructions allow it, but must
  not see two visualizations side by side. For Set 1, use the Scarf Plot for
  Question 1 and the video with fixation marker but without the K-Timeline for
  Question 2. For Set 3, use one view at a time and let the participant switch
  between the named views.
  \item \textbf{Two-view condition:} show the named complementary views at the
  same time. For Set 2, show Scarf Plot and AOI Timeline together for Question
  1, and synchronized video plus K-Timeline for Question 2. For Set 4, show the
  named views together as instructed.
\end{itemize}

The current application must have a true one-view study layout before this
booklet is used. If it only displays two panels, do not call that the one-view
condition. Keep image size, scale, selected records, AOI boundaries, video
episode, and playback position the same between matched groups.

\section*{Set 1 - Novice users, one view at a time}
\textbf{For the study leader:} Use the locked image stimulus and P1/P2 records.
Show one visualization at a time. For Question 1, display the Scarf Plot. For
Question 2, display the video and fixation marker without the K-Timeline.
Questions 3 and 4 are shared anchor questions; follow the one-view condition
while the participant answers. Allow about 25 minutes for the four questions.

\tasktitle{Question 1 - Visits to the Target AOI (H1)}
Compare P1 and P2 on the Scarf Plot. Which record has more visits to the Target
AOI? Use the visit definition in the study leader's introduction.
\choices{}{P1 has more visits.}{P2 has more visits.}{P1 and P2 have the same number of visits.}{A required record or view is missing.}

\tasktitle{Question 2 - Video episode (H1; linked to original H3)}
Watch episodes A and B with the fixation marker. Which episode shows primarily
focused inspection (focal attention)?
\choices{}{Episode A only.}{Episode B only.}{Both episodes.}{Neither episode.}

\tasktitle{Question 3 - Target-AOI dwell time (H2 shared question)}
Compare P1 and P2. Which record has the greater total dwell time in the Target
AOI? Use the sum of fixation durations inside the AOI.
\choices{}{P1 has the greater total dwell time.}{P2 has the greater total dwell time.}{P1 and P2 have equal total dwell time.}{A required record or view is missing.}

\tasktitle{Question 4 - Target-AOI dwell share (H2 shared question)}
Compare P3 and P4. Which record has the greater percentage of its total
fixation duration in the Target AOI?
\choices{}{P3 has the greater dwell share.}{P4 has the greater dwell share.}{P3 and P4 have equal dwell share.}{A required record or view is missing.}

\clearpage
\section*{Set 2 - Novice users, two views at the same time}
\textbf{For the study leader:} Use the same data and record mapping as Set 1.
Show the Scarf Plot and AOI Timeline together for Question 1. Show the video
and synchronized K-Timeline together for Question 2. Keep the shared anchor
questions and selected records the same as Set 1. Allow about 25 minutes for
the four questions.

\tasktitle{Question 1 - Visits to the Target AOI (H1)}
Compare P1 and P2 using the Scarf Plot and AOI Timeline together. Which record
has more visits to the Target AOI? A visit is one uninterrupted run of
fixations in that AOI.
\choices{}{P1 has more visits.}{P2 has more visits.}{P1 and P2 have the same number of visits.}{A required record or view is missing.}

\tasktitle{Question 2 - Video episode (H1; linked to original H3)}
Watch episodes A and B with the synchronized K-Timeline. Which episode shows
primarily focused inspection (focal attention)?
\choices{}{Episode A only.}{Episode B only.}{Both episodes.}{Neither episode.}

\tasktitle{Question 3 - Target-AOI dwell time (H2 shared question)}
Compare P1 and P2. Which record has the greater total dwell time in the Target
AOI? Use the sum of fixation durations inside the AOI.
\choices{}{P1 has the greater total dwell time.}{P2 has the greater total dwell time.}{P1 and P2 have equal total dwell time.}{A required record or view is missing.}

\tasktitle{Question 4 - Target-AOI dwell share (H2 shared question)}
Compare P3 and P4. Which record has the greater percentage of its total
fixation duration in the Target AOI?
\choices{}{P3 has the greater dwell share.}{P4 has the greater dwell share.}{P3 and P4 have equal dwell share.}{A required record or view is missing.}

\clearpage
\section*{Set 3 - Expert users, one view at a time}
\textbf{For the study leader:} Use the same locked study files and P1--P4
mapping used in all other sets. For Questions 1 and 2, participants may switch
between the named views, but show only one view at a time. Use the same episode
and record settings in both expert sets. Allow about 25 minutes for the four
questions.

\tasktitle{Question 1 - Target-AOI dwell share (H1)}
Compare P1 and P2. Which record has the greater percentage of its total
fixation duration in the Target AOI? Calculate dwell share as Target-AOI dwell
time divided by total fixation duration.
\choices{}{P1 has the greater dwell share.}{P2 has the greater dwell share.}{P1 and P2 have equal dwell share.}{A required record or view is missing.}

\tasktitle{Question 2 - Directed AOI transitions (H1)}
Compare P3 and P4. Which record has more transitions from AOI X directly to
AOI Y? Count each X-to-Y change once; do not count other transitions.
\choices{}{P3 has more X-to-Y transitions.}{P4 has more X-to-Y transitions.}{P3 and P4 have the same number of X-to-Y transitions.}{A required record or view is missing.}

\tasktitle{Question 3 - Target-AOI dwell time (H2 shared question)}
Compare the common anchor records P1 and P2. Which has the greater total dwell
time in the Target AOI?
\choices{}{P1 has the greater total dwell time.}{P2 has the greater total dwell time.}{P1 and P2 have equal total dwell time.}{A required record or view is missing.}

\tasktitle{Question 4 - Target-AOI dwell share (H2 shared question)}
Compare the common anchor records P3 and P4. Which has the greater percentage
of its total fixation duration in the Target AOI?
\choices{}{P3 has the greater dwell share.}{P4 has the greater dwell share.}{P3 and P4 have equal dwell share.}{A required record or view is missing.}

\clearpage
\section*{Set 4 - Expert users, two views at the same time}
\textbf{For the study leader:} Use the same locked study files and P1--P4
mapping used in all other sets. For Question 1, display the AOI Timeline and Attention Map together.
For Question 2, display the Transition Matrix and Scarf Plot together. Keep the
shared anchor questions and selected records the same as Set 3. Allow about 25
minutes for the four questions.

\tasktitle{Question 1 - Target-AOI dwell share (H1)}
Compare P1 and P2 using the AOI Timeline and Attention Map together. Which
record has the greater percentage of its total fixation duration in the Target
AOI? Calculate dwell share as Target-AOI dwell time divided by total fixation
duration.
\choices{}{P1 has the greater dwell share.}{P2 has the greater dwell share.}{P1 and P2 have equal dwell share.}{A required record or view is missing.}

\tasktitle{Question 2 - Directed AOI transitions (H1)}
Compare P3 and P4 using the Transition Matrix and Scarf Plot together. Which
record has more transitions from AOI X directly to AOI Y? Count each X-to-Y
change once; do not count other transitions.
\choices{}{P3 has more X-to-Y transitions.}{P4 has more X-to-Y transitions.}{P3 and P4 have the same number of X-to-Y transitions.}{A required record or view is missing.}

\tasktitle{Question 3 - Target-AOI dwell time (H2 shared question)}
Compare the common anchor records P1 and P2. Which has the greater total dwell
time in the Target AOI?
\choices{}{P1 has the greater total dwell time.}{P2 has the greater total dwell time.}{P1 and P2 have equal total dwell time.}{A required record or view is missing.}

\tasktitle{Question 4 - Target-AOI dwell share (H2 shared question)}
Compare the common anchor records P3 and P4. Which has the greater percentage
of its total fixation duration in the Target AOI?
\choices{}{P3 has the greater dwell share.}{P4 has the greater dwell share.}{P3 and P4 have equal dwell share.}{A required record or view is missing.}

\clearpage
\section*{Operator answer key and scoring sheet}
Complete this sheet from the locked source data before the first session. The
same answers must apply to every participant assigned to a matched set. Do not
infer answers from the response frequency or from the study hypothesis.

\begin{center}
\begin{tabular}{p{0.28\linewidth}p{0.32\linewidth}p{0.25\linewidth}}
\toprule
Question & Exact scoring rule & Correct option \\
\midrule
Set 1/2 Q1 & Compare visit counts: P1 vs P2 & \fillline{2cm} \\
Set 1/2 Q2 & Episode A/B focal duration share $>50\%$ & \fillline{2cm} \\
All sets Q3 & Compare Target-AOI dwell time: P1 vs P2 & \fillline{2cm} \\
All sets Q4 & Compare Target-AOI dwell share: P3 vs P4 & \fillline{2cm} \\
Set 3/4 Q1 & Compare Target-AOI dwell share: P1 vs P2 & \fillline{2cm} \\
Set 3/4 Q2 & Compare direct X-to-Y transition counts: P3 vs P4 & \fillline{2cm} \\
\bottomrule
\end{tabular}
\end{center}

For every participant, record the assigned set, experience group, answer to
each question, correctness (0 or 1), and time per question. Do not put the
answer key on participant copies. Report question-level and total accuracy;
compare matched layouts within each experience group for H1; compare novices
and experts on Questions 3 and 4 for H2.

\section*{Stop rule}
Do not administer a question if the study leader cannot identify one correct
option from the frozen data. Check especially that the two records are not
exactly tied unless the tie option is correct, all Target AOIs are visible, and
video episodes contain valid K-Coefficient values. Record a technical problem
separately and do not substitute another record during a session.

\end{document}
